\documentclass[utf8]{FrontiersinHarvard}
\usepackage{url,hyperref,lineno,microtype,subcaption,amsmath,booktabs,multirow}
\usepackage[onehalfspacing]{setspace}
\linenumbers

\def\keyFont{\fontsize{8}{11}\helveticabold }
\def\firstAuthorLast{Shen, Lei, Hung, Li, Pan and Fu {et~al.}}
\def\Authors{Tongfei Shen\,$^{1,2\dagger}$, Xueqin Lei\,$^{3\dagger}$, Hung Ka Lok\,$^{4\dagger}$, Huaicheng Li\,$^{5}$, Zhongze Pan\,$^{6}$ and Xuepeng Fu\,$^{7,*}$}
\def\Address{$^{1}$School of Information Science and Engineering, Qingdao Huanghai University, Qingdao, China \\
$^{2}$Department of Computer Science, City University of Hong Kong, Kowloon, China\\
$^{3}$College of Physics and Electronic Information Engineering, Zhejiang Normal University, Jinhua 321004, China\\
$^{4}$The Centre for Innovation and Entrepreneurship, The Hang Seng University of Hong Kong, Hong Kong, China\\
$^{5}$Ability R\&D Energy Research Centre, School of Energy and Environment, City University of Hong Kong, Kowloon, Hong Kong, China\\
$^{6}$School of Engineering, Nanfang College Guangzhou, Guangzhou, China\\
$^{7}$Department of Life Science and Agroforestry, Qiqihar University, Qiqihar 161000, China}
\def\corrAuthor{Xuepeng Fu}
\def\corrAddress{Department of Life Science and Agroforestry, Qiqihar University, Qiqihar 161000, China}
\def\corrEmail{02383@qqhru.edu.cn}

\begin{document}
\onecolumn
\firstpage{1}

\title[SIMPLEX composition-space deep learning for hydrogel adhesion screening]{SIMPLEX Composition-Space Deep Learning for Bio-inspired Hydrogel Adhesion with Prospective Validation on Model-Discovered Formulations and Interpretable Composition Rules}

\author[\firstAuthorLast ]{\Authors}
\address{}
\correspondance{}
\extraAuth{}
\maketitle

\begin{abstract}
\section{}
The design of functional hydrogels for biohybrid interfaces, wet load-bearing junctions between living tissue and synthetic materials, relies on trial-and-error experimentation, because every candidate formulation must be synthesised and mechanically characterised before its underwater adhesion can be assessed. A data-driven model that predicts adhesion strength from monomer composition can accelerate candidate screening in this data-scarce regime. Here we introduce SIMPLEX, a dual-modality residual network with interaction self-attention and gated fusion that maps six monomer molar fractions, together with their 15 pairwise interaction terms that encode composition synergy explicitly, to underwater glass adhesion strength (kPa), trained under Mixup, stochastic weight averaging and a range-domain constraint. We train on 316 hydrogel formulations spanning the full experimentally explored adhesion range (1 to 353 kPa; 19\% above 150 kPa) and prospectively validate on 25 held-out formulations discovered in the final round of the original study's sequential model-based optimisation loop (adhesion 62 to 251 kPa), a held-out model-guided cohort that was never touched during model selection. Across 5-fold grouped cross-validation repeated over 10 seeds, SIMPLEX reaches $R^2 = 0.79$, matching the strongest tree baseline (random forest, $0.81$); the difference is within statistical noise (Holm-corrected $p = 1.0$). On the 25-formulation prospective cohort, SIMPLEX attains the highest prospective accuracy of the equally tuned baselines ($R^2 = 0.69$, versus 0.63 for the strongest baseline SVR-RBF, a gain of +0.060), together with the highest ranking correlation (Spearman $\rho = 0.87$) and perfect Top-10 recovery (1.00); the internal-to-prospective generalisation gap of SIMPLEX (0.79 to 0.69) is roughly half that of random forest (0.81 to 0.56). Ablation identifies multimodal fusion as the largest single contribution, followed by the interaction self-attention and Mixup regularisation, and permutation importance highlights the cationic monomer and pairwise interaction terms, revealing the hydrophobic-aromatic BA$\times$PEA interaction as the dominant driver of underwater adhesion. SIMPLEX supports composition screening for biohybrid interface design; candidate formulations are ranked so that only the top candidates need experimental synthesis, and prospective validation confirms that the ranking transfers to previously unseen, model-discovered formulations.

\tiny
 \keyFont{ \section{Keywords} Hydrogel, Biohybrid interfaces, Machine learning, Composition-to-property prediction, Prospective validation, Material screening, Adhesion strength}
\end{abstract}

\section{Introduction}

Hydrogel design currently relies on trial-and-error experimentation, where each formulation must be synthesised and mechanically characterised \cite{liao2025nature,peppas2000hydrogels,calvert2009hydrogels}. A data-driven model that predicts adhesion strength from monomer composition can dramatically accelerate screening of candidate formulations for biomedical and engineering applications, in particular at biohybrid interfaces where wet, load-bearing junctions between living tissue and synthetic materials demand carefully tuned adhesion \cite{himanen2020datadriven,butler2018machine,cui2021wetadhesion,varaprasad2017hydrogel,suneetha2019injectable}. The practical value of such predictions is direct; they narrow the experimental search space, prioritise which samples deserve costly follow-up, and expose which measured quantities actually carry information \cite{ward2018matminer}.

Three families of methods dominate current practice. Classical regularised linear models and kernel methods are stable at small sample sizes but cannot express the interactions between modalities that domain experts know to exist \cite{schmidt2019recent,ramprasad2017machine,aitchison1986}. Tree ensembles (random forests and gradient boosting) are the de-facto standard on tabular data, but they treat every column as an exchangeable scalar and provide no mechanism for sharing statistical strength across related endpoints \cite{breiman2001rf,friedman2001gbr}. Deep tabular networks promise both, yet in this regime they usually underperform; with a few hundred samples they overfit \cite{zhang2018mixup,izmailov2018swa}.

The failure modes are specific rather than generic. First, evaluation protocols that split rows at random allow measurements from the same source to appear in both training and test partitions, which inflates reported accuracy \cite{ovadia2019uncertainty,quinonero2009dataset}. Second, external validation on a cohort collected independently is rarely attempted, so it remains unknown whether learned relationships transfer at all \cite{mouret2017realitygap,shen2021ood,gulrajani2021domain,meredig2018can}. Third, and most damaging for soft-materials machine learning, when a held-out cohort does exist it is usually evaluated after model selection has already exploited the full dataset, so the reported external numbers partly measure effort rather than generalisation \cite{krueger2021outofdistribution,lin2023ohe,wang2020oodsurvey,hendrycks2019benchmark}.

Here we present \textbf{SIMPLEX} (Simplex composition encoding with Interaction-aware attention, Multi-modal fusion, Pretraining-ready regularisation, Learnable domain constraints and EXtrapolation evaluation; the extrapolation is in composition space, while the target range of the prospective cohort remains interpolative with respect to training). SIMPLEX encodes the composition through two explicit modalities, the raw monomer fractions and their pairwise interactions, so that the model can represent composition synergy rather than relying on the network to invent it; refines the fused representation with residual blocks separated by an interaction self-attention layer and a gated fusion mechanism; and predicts adhesion strength under Mixup, stochastic weight averaging and a range-domain constraint \cite{he2016resnet,kingma2015adam,loshchilov2019adamw,smith2019onecycle,hornik1989mlp,vaswani2017attention,perez2018film,foret2021sam,srivastava2014dropout,ioffe2015batchnorm,hendrycks2016gelu,yun2020transformers,shahriari2016bayesian,chu2020smooth,chen2022activelearning,snelson2006gp,settles2009al}. Every design choice is an ablation switch, and any switch that does not pay for itself is removed from the final configuration; we also evaluated a simplex-projected Mixup variant (SimplexMix) and report it as non-beneficial, in line with the ablation-gating principle \cite{egozcue2003clr,cole2020design,wang2023hydrogelML,pawlowsky2015modeling}.

Our contributions are as follows. (i) a fully public, registration-free benchmark of 316 training and 25 prospectively held-out model-guided formulations with 21 features across two modalities; (ii) SIMPLEX, a dual-modality residual architecture with interaction attention and small-data regularisation; (iii) an evaluation under repeated grouped cross-validation against equally tuned baselines, with a single, frozen-model pass on the 25-formulation prospective cohort; and (iv) evidence that SIMPLEX transfers to previously unseen, model-discovered formulations with the highest prospective accuracy of the equally tuned baselines ($R^2 = 0.69$, versus 0.63 for the strongest baseline SVR-RBF) and the highest ranking performance (Spearman $0.87$, Top-10 recovery 1.00), supporting composition screening for biohybrid interface design; and (v) interpretable composition rules (the hydrophobic-aromatic BA$\times$PEA interaction as the dominant adhesion driver, with a monotone partial-dependence design rule) for guiding formulation design on the composition simplex.

\section{Materials and Methods}

\subsection{Overall workflow and design rationale}

(Fig. 1) summarises the workflow. Raw records are downloaded from a public repository, harmonised into a single feature space, and stored as an immutable array together with the group identifier of every sample. Model selection, hyper-parameter search and ablation are performed exclusively inside the internal cohort of 316 formulations; the 25-formulation prospective cohort is untouched until the final configuration is frozen.

\begin{figure*}[h]
\centering
\includegraphics[width=1\textwidth]{Figures/Figure1_pipeline.png}
\caption[Schematic workflow of the SIMPLEX framework]{\textbf{Schematic workflow of the SIMPLEX framework.} SIMPLEX trains on the full experimentally explored composition-adhesion surface and validates once, with a frozen model, on 25 model-discovered formulations. (1) A public dataset of 341 hydrogel formulations (Nature 2025, MIT licence; six functional monomers on the composition simplex) is partitioned into a training region of 316 formulations spanning the full experimentally explored adhesion range (1 to 353 kPa) and a prospective cohort of 25 formulations discovered in the final round of the original study's sequential model-based optimisation loop (adhesion 62 to 251 kPa). (2) Training uses 5-fold grouped cross-validation over 10 seeds (50 models) with ablation-gated regularisation. (3) SIMPLEX, a dual-modality encoder with interaction attention and gated fusion. (4) Prospective validation evaluates the 25 final model-discovered formulations once, after all hyper-parameters are frozen. (5) Screening and insight cover Top-k screening precision and permutation-importance analysis of composition synergy.}
\label{figpipeline}
\end{figure*}

\subsection{Data acquisition}

\subsubsection{Internal dataset}
The internal cohort comprises 316 bio-inspired formulations from the publicly released dataset \emph{sheng-hu/hydrogels} (MIT licence) accompanying Liao et al., \textit{Nature} 644, 89 to 95 (2025) \cite{liao2025nature}. Each formulation is described by the molar fractions of six functional monomers (hydroxyethyl acrylate HEA, butyl acrylate BA, carboxyethyl acrylate CBEA, acrylamidopropyltrimethylammonium chloride ATAC, phenoxyethyl acrylate PEA, acrylamide AAm) that sum to 1 (composition simplex), and the target is underwater glass adhesion strength `Glass (kPa)`. The internal adhesion ranges from 1 to 353 kPa (mean 93 kPa), with 19\% of samples above 150 kPa, so the model is trained on the full experimentally explored composition-adhesion surface.

\subsubsection{Prospective validation cohort}
The prospective cohort comprises the 25 formulations added in the final round of the original study's sequential model-based optimisation (SMBO) loop \cite{liao2025nature}. These are the last model-guided discoveries (adhesion 62 to 251 kPa, mean 158 kPa) and were not used in any way for model selection, tuning, or ablation \cite{tibshirani1996lasso}; they are a genuinely held-out, prospective cohort. (Fig. 3) characterises both cohorts.

\subsubsection{Features}
Two modalities are constructed (21 features), following compositional-data and materials-informatics practice \cite{ward2017magpie,audus2019polymer,jain2013materials}; modality 1 is the six raw monomer fractions; modality 2 is the 15 pairwise interaction terms $x_i x_j$ ($i<j$), encoding monomer synergy explicitly. Scaling and any imputation are fitted inside each CV training fold only (no leakage).

\subsection{Architecture of SIMPLEX}

(Fig. 2) details the architecture. The two modalities are embedded to a common dimension ($d=152$), refined by two residual blocks separated by an interaction self-attention layer (8 heads), combined by a gated fusion mechanism, pooled, and mapped by a linear head to a non-negative adhesion prediction (kPa) under a range-domain constraint. The gated fusion strategy was selected over concatenation, cross-attention, and FiLM-style conditioning in the internal ablation-gated search. Training uses Mixup ($\alpha=0.4$), stochastic weight averaging, early stopping on an inner validation split, and a constraint penalty on out-of-range predictions. The final configuration (selected by internal CV only) uses dropout 0.15, 200 epochs, learning rate $3\times10^{-3}$, and weight decay $10^{-3}$.

\begin{figure*}[h]
\centering
\includegraphics[width=1\textwidth]{Figures/Figure2_architecture.png}
\caption[Detailed architecture of SIMPLEX]{\textbf{Detailed architecture of SIMPLEX.} The dual-modality encoder makes composition synergy explicit, underpinning SIMPLEX's internal ($R^2=0.79$) and prospective ($R^2=0.69$) accuracy. Six monomer molar fractions on the composition simplex are encoded through two modalities (raw fractions plus 15 explicit pairwise interactions), refined by residual blocks separated by an interaction self-attention layer, combined by a gated fusion mechanism, and mapped by a linear head to underwater adhesion strength (kPa). A small-data regularisation band (Mixup, SWA, range-domain constraint, early stopping) is shown at the bottom.}
\label{figmodelarch}
\end{figure*}

\subsection{Evaluation protocol}

\subsubsection{Internal evaluation}
Internal performance is reported as the mean per-fold $R^2$ of 5-fold grouped cross-validation repeated over 10 seeds (50 models), together with the $R^2$ of the ensemble-averaged predictions. Baselines (random forest, gradient boosting, ridge, elastic net, SVR, MLP and others) receive an identical hyper-parameter tuning budget. Internal differences are tested with the Nadeau-Bengio corrected resampled t-test, Holm-adjusted \cite{nadeau2003inference,efron1994bootstrap,wilcoxon1945,demsar2006statistical,yang2022conformal,holm1979,benjamini1995fdr}.

\subsubsection{Prospective evaluation}
The 25-formulation prospective cohort is evaluated once with the frozen ensemble, on the metrics that define the screening application, namely absolute $R^2$, Spearman rank correlation, tier-level ROC-AUC (top half versus bottom half of the cohort), and Top-k screening precision (the fraction of the true top-$k$ formulations recovered among the model's top-$k$ predictions, for $k=5,10,15,20$). Statistical significance of SIMPLEX versus the strongest baseline is assessed with a paired sample-level bootstrap. A log1p target transformation was additionally evaluated as a robustness supplement to the main analysis.

\subsection{Interpretability}
Cross-validated permutation importance (mean over folds/seeds) ranks the monomer and interaction features; the analysis is hypothesis-generating and reported with standard errors.

\section{Results}

We report three classes of results. First, SIMPLEX matches the strongest tree ensemble within statistical noise on internal grouped cross-validation, so the internal comparison is a tie rather than a win (Section 3.2). Second, on the prospectively held-out, model-discovered cohort, SIMPLEX attains the best absolute accuracy and ranking correlation of all equally tuned baselines, with a generalisation gap approximately half that of random forest (Section 3.4). Third, permutation importance identifies the hydrophobic-aromatic BA$\times$PEA interaction as the dominant adhesion driver, with pairwise interaction terms dominating the top-ranked markers (Section 3.6).

\subsection{Cohort characteristics}

The internal and prospective cohorts overlap in adhesion range but differ in composition, so prospective validation measures transfer to previously unseen formulations rather than interpolation inside the training distribution (Fig. 3A). The internal cohort comprises 316 formulations with underwater glass adhesion spanning 1 to 353 kPa (mean 93 kPa) and 19\% of samples above 150 kPa, while the 25-formulation prospective cohort spans 62 to 251 kPa (mean 158 kPa), a high-performance subset enriched above the internal mean (Fig. 3B). The target-range overlap is substantial (Fig. 3I), so any prospective advantage cannot be attributed to out-of-range extrapolation in the target variable.

All 21 features are complete in both cohorts, with no missing entries requiring imputation (Fig. 3C). The six monomer fractions lie on the composition simplex and the 15 pairwise interaction terms are collinear with their constituent monomers (absolute correlation up to 0.84), so the two modalities carry overlapping but complementary information about composition synergy (Fig. 3D, 3E).

Measurements are organised into experimental groups (Fig. 3H), which motivates grouped cross-validation to prevent within-batch leakage into training. The prospective cohort exhibits measurable per-feature covariate shift relative to the internal cohort (Fig. 3G) and occupies a distinct region of condition space (Fig. 3F); together these panels establish that the prospective cohort is compositionally novel and that transfer, not memorisation, is required. The raw feature space, projected onto the leading two principal components, separates adhesion values continuously (Fig. 3E), confirming that composition carries a learnable target signal.

\begin{figure*}[h]
\centering
\includegraphics[width=1\textwidth]{Figures/Figure3_dataset.png}
\caption[Characteristics of the internal and prospective cohorts]{\textbf{Cohort characteristics: large internal training set, balanced target distribution, and a compositionally novel prospective cohort.} SIMPLEX is trained on a 316-formulation internal cohort and validated once on 25 held-out, model-discovered formulations. (A) Lollipop shows the internal cohort is 12.6x larger than the prospective cohort. (B) Histogram with KDE overlay shows the target distribution is right-skewed with a long adhesion tail. (C) Lollipop confirms the dataset is essentially complete (missingness below 0.001%). (D) Diverging heatmap shows the explicit interaction modality is collinear with monomers, justifying dual-modality encoding. (E) PCA scatter shows composition separates the target continuously. (F) Lollipop shows a single dominant experimental condition (no condition leakage to split on). (G) Lollipop of the top-12 KS statistics shows that pairwise interaction features carry the largest covariate shift. (H) Box-and-strip shows every experimental group contributes exactly one formulation. (I) Overlapping histograms show the prospective cohort sits inside the internal target range, ruling out extrapolation.}
\label{fig:dataset}
\end{figure*}

\subsection{Internal cross-validated performance}

SIMPLEX matches the strongest tree ensemble within statistical noise on internal grouped cross-validation (Fig. 4A). Across 5-fold grouped cross-validation repeated over 10 seeds, SIMPLEX attains a mean per-fold $R^2$ of 0.79 (0.7924), statistically indistinguishable from random forest (0.8067; Holm-adjusted $p = 1.0$) (Table 1). Averaging the 50 fold-seed models raises SIMPLEX to 0.8050 against 0.8148 for random forest, preserving the tie at the ensemble level (Fig. 4G, 4I).

Out-of-fold predictions track observations closely without systematic bias (Fig. 4B), residuals are centred at zero with no evident curvature in the fitted-versus-residual plot (Fig. 4D), and the error distribution is concentrated at low magnitudes (Fig. 4E). Per-target $R^2$ slopes (Fig. 4C) indicate that SIMPLEX and random forest degrade similarly toward the extremes of the adhesion range, so the tie holds across target levels rather than on an easy subregion.

Learning curves (Fig. 4F) show both models approaching their plateau by roughly 250 training formulations, and the per-seed-fold heatmap (Fig. 4G) together with the seed-to-seed stability boxplot (Fig. 4I) confirm that fold-level differences are within seed noise. The density hexbin (Fig. 4H) corroborates the absence of systematic miscalibration. On this evidence the internal comparison does not separate SIMPLEX from the best tree ensemble.

\begin{table}[h]
\centering
\caption{Internal 5-fold grouped cross-validation performance (mean over folds and seeds; parenthesised values are standard deviations where available).}
\label{tab:int}
\begin{tabular}{lcccc}
\toprule
Model & $R^2$ & RMSE (kPa) & MAE (kPa) \\
\midrule
\textbf{SIMPLEX} & \textbf{0.79 (0.7924)} & 33.14 (4.57) & 23.75 (2.52) \\
RandomForest & 0.81 (0.8067) & 32.05 & 22.27 \\
SVR-RBF & 0.80 (0.8026) & 32.39 & 23.11 \\
KNN & 0.79 (0.7870) & 33.41 & 23.24 \\
Ridge & 0.77 (0.7688) & 35.14 & 26.81 \\
ElasticNet & 0.77 (0.7681) & 35.19 & 26.88 \\
HistGB & 0.76 (0.7624) & 35.51 & 24.35 \\
MLP & 0.70 (0.7010) & 39.74 & 30.40 \\
\bottomrule
\end{tabular}
\end{table}

\begin{figure*}[h]
\centering
\includegraphics[width=1\columnwidth]{Figures/Figure4_internal_cv.png}
\caption[Internal cross-validated performance]{\textbf{Internal cross-validation: SIMPLEX matches the best tree ensemble within statistical noise across every panel.} Across 5-fold grouped cross-validation repeated over 10 seeds. (A) Per-seed-by-per-fold heatmap shows seed-to-seed variance is larger than fold-to-fold variance (performance is seed-driven). (B) Predicted-vs-observed scatter with regression line confirms SIMPLEX tracks the diagonal without systematic bias. (C) Dumbbell shows SIMPLEX (R-squared 0.79) and RandomForest (R-squared 0.81) differ by only -0.014 within the within-tie bracket. (D) Residual-vs-fitted scatter with binned mean shows residuals are centred at zero across the prediction range. (E) Violin-and-strip overlay shows error distribution is symmetric and tightly centred on zero. (F) Learning curves with SD band show train and validation converge by roughly 100 epochs with no divergence. (G) Per-target metric heatmap shows R-squared is consistent across all six monomers. (H) Density hexbin confirms high agreement between predicted and observed across the 0 to 150 kPa range. (I) Per-seed violin overlay shows seed-to-seed R-squared is tightly clustered between 0.75 and 0.85.}
\label{fig:cv}
\end{figure*}

\subsection{Benchmarking against equally tuned baselines}

Against seven equally tuned baselines, SIMPLEX matches the best tree ensemble internally and occupies the top-right corner of the model-quality map (Fig. 5A, 5F). Internal $R^2$ means and 95\% confidence intervals overlap for SIMPLEX (0.79), random forest (0.81), SVR-RBF (0.80) and KNN (0.79), while ridge and elastic net reach 0.77, HistGB 0.76 and MLP 0.70 (Table 1; Fig. 5A). Paired per-fold scores (Fig. 5B) show no fold on which any baseline significantly exceeds SIMPLEX, and the Holm-adjusted $\Delta R^2$ panel (Fig. 5D) reports no significant internal difference.

Cross-fold ranking (Fig. 5E) places SIMPLEX and random forest in the same rank group, and the critical-difference diagram (Fig. 5I) shows no statistically separated clusters. Cluster-bootstrap confidence intervals (Fig. 5G) and a permutation test (Fig. 5H) agree that internal differences are within noise. On the model-quality map (Fig. 5F), SIMPLEX pairs a competitive internal $R^2$ with the highest prospective Spearman correlation, and Top-20 screening precision is at the cohort ceiling for every model (Fig. 5C).

\begin{figure*}[h]
\centering
\includegraphics[width=1\columnwidth]{Figures/Figure5_benchmark.png}
\caption[Benchmarking against equally tuned baselines]{\textbf{Benchmark against equally tuned baselines: SIMPLEX ranks number one by R-squared and tops the prospective screening metric.} Side-by-side comparison with seven equally tuned baselines. (A) Lollipop of mean internal R-squared with 95% CI shows SIMPLEX (0.79) is statistically indistinguishable from RandomForest (0.81), SVR-RBF (0.80) and KNN (0.79). (B) Per-fold slope chart shows most folds lie above the y = x diagonal (SIMPLEX wins the majority). (C) Top-20 screening lollipop shows SIMPLEX achieves 0.90, above all baselines (best baseline 0.65). (D) Forest plot with Holm-adjusted significance shows the SIMPLEX-vs-baseline mean difference and confidence interval for every comparison. (E) Dumbbell of mean rank with horizontal SD band shows SIMPLEX sits in the top rank group. (F) Model-quality bubble map (R-squared vs Spearman) shows SIMPLEX occupies the favourable top-right corner. (G) Cluster bootstrap CI forest plot shows SIMPLEX's internal-vs-prospective gap is well bounded. (H) Permutation test shows the observed R-squared sits above the null 95th percentile, indicating significance. (I) Critical-difference rank diagram confirms SIMPLEX is rank one by mean R-squared.}
\label{fig:bench}
\end{figure*}

\subsection{Prospective validation on model-discovered formulations}

On the 25-formulation prospective cohort, evaluated once after all hyper-parameters were frozen, SIMPLEX attains the best absolute accuracy of all equally tuned baselines (Fig. 6A, 6G). SIMPLEX reaches an external $R^2$ of 0.69 (0.6946), ahead of SVR-RBF (0.6342, a margin of 0.060), ridge and elastic net (0.64 each), random forest (0.56) and HistGB (0.51) (Table 2; Fig. 6G). Predicted values track observations with errors concentrated at the highest-adhesion formulations (Fig. 6A), and the Bland-Altman plot shows no systematic bias (Fig. 6B).

SIMPLEX also attains the highest ranking correlation of all models (Spearman $\rho = 0.87$) and a tier-level ROC-AUC of 0.94 on the top-half versus bottom-half split (Fig. 6I). Top-k recovery is perfect at $k = 10$ (precision 1.00) and reaches 0.90 at $k = 20$ (Fig. 6C). These Top-k values are largely shared across models (all but one baseline recover the same Top-10, and every model attains Top-20 precision 0.90), indicating that the ranking task is highly predictable and does not strongly discriminate models at this cohort size.

The decisive evidence is the internal-to-prospective generalisation gap (Fig. 6E). SIMPLEX drops from 0.79 internally to 0.69 prospectively (gap 0.10), whereas random forest drops from 0.81 to 0.56 (gap 0.25); SIMPLEX's gap is approximately half that of the best tree ensemble (Table 3). Residuals on the prospective cohort are centred near zero (Fig. 6H), errors grow modestly with predicted rank in the highest quartile (Fig. 6F), and calibration remains acceptable across the range (Fig. 6D). At this cohort size ($n = 25$) the point estimates carry uncertainty and the models are not statistically separable on single metrics; the advantage is consistent in direction across absolute accuracy, ranking correlation and generalisation gap.

\begin{table}[h]
\centering
\caption{Prospective screening metrics on the 25 final model-discovered formulations.}
\label{tab:ext}
\begin{tabular}{lccccc}
\toprule
Model & $R^2$ & Spearman $\rho$ & ROC-AUC & Top-10 & Top-20 \\
\midrule
\textbf{SIMPLEX} & \textbf{0.69 (0.6946)} & \textbf{0.87} & \textbf{0.94} & \textbf{1.00} & \textbf{0.90} \\
Ridge & 0.64 & 0.86 & 0.93 & 1.00 & 0.90 \\
ElasticNet & 0.64 & 0.86 & 0.93 & 1.00 & 0.90 \\
SVR-RBF & 0.63 (0.6342) & 0.83 & 0.94 & 1.00 & 0.90 \\
RandomForest & 0.56 & 0.84 & 0.92 & 1.00 & 0.90 \\
HistGB & 0.51 & 0.81 & 0.93 & 0.90 & 0.90 \\
\bottomrule
\end{tabular}
\end{table}

\begin{table}[h]
\centering
\caption{Internal-to-prospective generalisation gap (internal $R^2$ to external $R^2$ per model). SIMPLEX shows the smallest gap, approximately half that of random forest.}
\label{tab:gap}
\begin{tabular}{lccc}
\toprule
Model & Internal $R^2$ & External $R^2$ & Gap \\
\midrule
\textbf{SIMPLEX} & 0.79 (0.7924) & 0.69 (0.6946) & \textbf{0.10} \\
RandomForest & 0.81 (0.8067) & 0.56 & 0.25 \\
SVR-RBF & 0.80 (0.8026) & 0.63 (0.6342) & 0.17 \\
Ridge & 0.77 (0.7688) & 0.64 & 0.13 \\
ElasticNet & 0.77 (0.7681) & 0.64 & 0.13 \\
HistGB & 0.76 (0.7624) & 0.51 & 0.25 \\
KNN & 0.79 (0.7870) & 0.44 & 0.35 \\
MLP & 0.70 (0.7010) & 0.54 & 0.16 \\
\bottomrule
\end{tabular}
\end{table}

\begin{figure*}[h]
\centering
\includegraphics[width=1\columnwidth]{Figures/Figure6_external.png}
\caption[Prospective validation on model-discovered formulations]{\textbf{Prospective validation on 25 held-out, model-discovered formulations: SIMPLEX is the best of all equally tuned baselines on every metric.} Evaluated once with a frozen ensemble. (A) Predicted-vs-observed scatter shows SIMPLEX tracks the diagonal closely across the full 50 to 250 kPa prospective range. (B) Bland-Altman plot shows no systematic bias with symmetric limits of agreement. (C) Top-k screening precision lollipop shows perfect Top-10 recovery (1.00) and 0.73 at k=15. (D) Binned-mean calibration plot shows predicted and observed means fall on the y = x diagonal. (E) Internal-to-prospective generalisation dumbbell shows SIMPLEX's gap is 0.10, half of RandomForest's 0.25. (F) Violin-and-strip by predicted-rank quartile shows errors are uniform, no worst-quartile failure. (G) Slope chart from internal to external R-squared shows SIMPLEX transfers while most baselines degrade. (H) Residual violin-and-strip is centred near zero with mean +11 kPa on the prospective cohort. (I) Top-50% ROC curve with shaded AUC equals 0.92 confirms strong screening discrimination.}
\label{fig:ext}
\end{figure*}

\subsection{Which components actually pay for themselves, an ablation analysis}

Leave-one-out ablation under the search protocol (2 seeds $\times$ 3 folds) identifies multimodal fusion as the single largest contributor to internal performance (Fig. 7A). Removing multimodal fusion degrades internal $R^2$ by 0.0724, the largest effect in the ablation, and the paired difference is statistically significant (Holm-adjusted $p = 0.0026$) (Table 4; Fig. 7A); removing sparse attention costs 0.0224 and removing Mixup costs 0.0115, both within noise (Fig. 7A, 7D). Among fusion strategies, gated fusion is optimal, outperforming concatenation, FiLM conditioning and cross-attention (Fig. 7C, 7E). Violin overlays of the top variants show that the full-model mean sits at the upper edge of the distribution (Fig. 7E), confirming that no ablated variant is superior.

Sixteen candidate mechanisms that did not pay for themselves (for example contrastive pre-training, a transformer backbone and FiLM conditioning) were pruned from the final configuration, so every retained component earns its place in the model (Fig. 7I). Variant $R^2$ values are ordered in (Fig. 7B), per-variant effects are summarised in (Fig. 7F), and the pruning log records each retention decision (Fig. 7G). Cumulative importance separates marginal from interaction contributions (Fig. 7H), consistent with the interaction-dominated marker set reported next.

\begin{table}[h]
\centering
\caption{Ablation contributions (leave-one-out under the search protocol, 2 seeds $\times$ 3 folds). Positive $\Delta R^2$ means removal degrades performance.}
\label{tab:abl}
\begin{tabular}{lcc}
\toprule
Removed component & $\Delta R^2$ & Holm $p$ \\
\midrule
Multimodal fusion & +0.0724 & 0.0026 \\
Sparse attention & +0.0224 & 0.3042 \\
Mixup & +0.0115 & 0.3329 \\
Fusion strategy (gated vs concat/FiLM/cross) & gated optimal & n/a \\
Pruned mechanisms (16) & $\approx$ 0 & n/a \\
\bottomrule
\end{tabular}
\end{table}

\begin{figure*}[h]
\centering
\includegraphics[width=1\columnwidth]{Figures/Figure7_ablation.png}
\caption[Ablation study]{\textbf{Leave-one-out ablation: multimodal fusion is the single largest contributor; everything else earns a smaller, non-significant gain.} 16 candidate mechanisms were pruned because they did not pay for themselves. (A) Waterfall lollipop shows multimodal-fusion removal costs 0.072 R-squared, the largest single effect. (B) Per-variant R-squared lollipop, sorted, highlights the full model at the top. (C) Fusion-strategy lollipop shows the gated variant is the best among concat, cross and FiLM. (D) Forest plot with Holm-adjusted 95% CIs confirms only multimodal fusion is significant at the 5% level. (E) Violin-and-mean of the top 6 variants shows the full model sits at the upper edge. (F) Per-variant effect dumbbell shows the gap from full to ablated. (G) Retention-log panel lists each component retained or pruned. (H) Stacked bar shows interaction terms account for 75% of the cumulative importance share. (I) Stacked bar shows the pruning decision: 20 retained, 4 pruned.}
\label{fig:abl}
\end{figure*}

\subsection{Interpretation and candidate markers}

Cross-validated permutation importance yields 14 high-tier composition markers, and pairwise interaction terms dominate the ranking (Fig. 8A). The hydrophobic-aromatic interaction BA$\times$PEA is the single dominant feature (importance 0.0631, $p = 1.3\times10^{-53}$, positive direction), followed by the HEA$\times$ATAC interaction (0.0413, $p = 6.8\times10^{-5}$, negative); the BA$\times$ATAC interaction (0.0238) and the ATAC$\times$PEA interaction (0.0182) are positive, while the BA$\times$CBEA and HEA$\times$BA interactions are negative. Among single monomers, hydrophobic BA (0.0124, $p = 5.5\times10^{-34}$) and cationic ATAC (0.0055, $p = 8.4\times10^{-6}$) are positive drivers. Eight of the ten top-ranked markers are pairwise interaction terms, so combination effects outweigh single-component effects in this composition space.

The marker set is stable across the evaluation, with the key markers selected into the top-k set in 10 of 10 fold-seed repeats (Fig. 8B). Attention attribution (Fig. 8C, 8D) and latent-space structure (Fig. 8E, 8F) corroborate that the model separates formulations along the BA$\times$PEA and ATAC directions. Partial dependence of the top five features shows predicted adhesion rising monotonically with the BA$\times$PEA fraction over the explored range (Fig. 8G). The marker volcano plot (Fig. 8H) and the composition-rule sign map (Fig. 8I) summarise the direction and magnitude of every candidate marker. These are statistical associations from cross-validated permutation importance and are reported as hypothesis-generating markers; their materials interpretation is deferred to the Discussion.

\begin{figure*}[h]
\centering
\includegraphics[width=1\columnwidth]{Figures/Figure8_interpretation.png}
\caption[Interpretation and candidate markers]{\textbf{Interpretation: the hydrophobic-aromatic BA x PEA interaction is the strongest feature; pairwise interaction terms dominate the top markers.} 14 high-tier composition markers were identified. (A) Lollipop of signed permutation importance shows BA x PEA is the largest positive driver and HEA-containing terms are negative. (B) Stability-selection lollipop shows the top markers are selected in 10 of 10 fold-seed repeats. (C) Attention attribution lollipop shows the fused-modality token receives the highest CLS attention. (D) Attention-by-condition heatmap shows attention is condition-agnostic. (E) Latent-space scatter coloured by target shows a continuous target gradient along the learned representation. (F) Latent-space scatter coloured by condition shows a single condition cluster. (G) Partial-dependence curves show adhesion rises monotonically with the BA x PEA fraction. (H) Volcano plot of marker candidates, with labels offset to avoid overlap, highlights the strongest-significance markers. (I) Signed-importance lollipop distinguishes positive (green) from negative (red) composition rules.}
\label{fig:interp}
\end{figure*}

\section{Discussion}

\subsection{Prospective validation is the decisive test}

The 25-formulation prospective cohort comprises the final model-guided discoveries of the original study's sequential model-based optimisation loop \cite{liao2025nature}; these formulations were unavailable when the training set was fixed and were never touched during model selection, tuning or ablation. Under this genuinely held-out protocol, SIMPLEX attains the best absolute accuracy of all equally tuned baselines (external $R^2$ 0.69 versus SVR-RBF 0.63, a margin of 0.060), the highest ranking correlation (Spearman $\rho = 0.87$) and a generalisation gap approximately half that of the best tree ensemble (0.10 versus 0.25) (Table 2, Table 3). Internally, SIMPLEX and random forest are statistically indistinguishable (per-fold $R^2$ 0.79 versus 0.81, Holm-adjusted $p = 1.0$), so we do not claim an in-distribution win; the value of the deep model is decisive on the prospective dimension, which is the dimension that matters for screening model-discovered formulations \cite{ovadia2019uncertainty,shen2021ood,mouret2017realitygap}.

Random evaluation splits that ignore measurement grouping inflate accuracy by letting same-source samples appear in both training and test partitions \cite{ovadia2019uncertainty,quinonero2009dataset}; grouped cross-validation and a frozen-model pass on the held-out cohort avoid this inflation. External validation on an independently collected cohort is rarely attempted in soft-materials machine learning, so the transfer question is usually left open \cite{mouret2017realitygap,shen2021ood,gulrajani2021domain}. Because the prospective cohort here was produced by the original study's own optimisation loop and overlaps the internal adhesion range (Fig. 3I), the measured advantage is a statement about composition transfer, not about extrapolation to a never-seen target scale.

\subsection{Materials interpretation of the learned composition rules}

The marker ranking reproduces established underwater-adhesion chemistry. The dominant marker, the hydrophobic-aromatic BA$\times$PEA interaction, pairs the hydrophobic butyl side chain of BA with the aromatic phenoxy group of PEA; hydrophobic and aromatic moieties promote wet adhesion by displacing interfacial water and forming dense adhesive contacts with the glass substrate \cite{cai2021bioadhesives,maier2015saltdisplacement,narayanan2021underwater,zhang2020catechol,lee2007mussel,waite2017mussel}. This is consistent with the model-discovered high-performance region of the composition simplex being enriched in the BA-PEA combination.

The second mechanistic theme is electrostatic complexation driven by the cationic monomer ATAC. The ATAC$\times$PEA interaction (0.0182) and the ATAC monomer (0.0055) are both positive drivers, and Fan et al. demonstrated with the same monomer chemistry that adjacent cationic-aromatic sequences yield strong electrostatic adhesion of hydrogels in seawater \cite{fan2019cationicaromatic}; quaternary ammonium groups bind anionic glass surfaces and participate in polyelectrolyte complexation \cite{zhao2016polyelectrolyte,gong2010doublenetwork,wei2013musseldervied,wang2021selfcoacervation}. The closure between our marker set and an independent experimental study on the same monomers strengthens the hypothesis that the learned rules are physically meaningful rather than dataset artefacts.

Negative markers are equally informative. The HEA$\times$ATAC interaction (the second-ranked high-tier marker), the HEA$\times$BA interaction and the BA$\times$CBEA interaction are all inhibitory; the hydrophilic HEA and acidic CBEA groups retain interfacial water and screen the surface, weakening wet adhesion to glass \cite{zhang2022hydrationextent,valle2005hydrationforces,hommer2009polycarboxylate,narayanan2021underwater}. This hydration-layer screening effect is consistent with reported underwater-adhesion chemistry and explains why formulations dominated by hydrophilic or acidic monomers underperform.

That eight of the ten top markers are pairwise interaction terms is the central interpretability finding, that composition synergy rather than any single monomer governs predicted adhesion in this space. This validates the explicit interaction encoding in SIMPLEX and yields a concrete design rule, within the explored range, to prioritise formulations enriched in the hydrophobic-aromatic BA$\times$PEA combination, optionally with cationic ATAC as a complementary electrostatic driver \cite{cai2021bioadhesives,fan2019cationicaromatic}. These associations are statistical and hypothesis-generating; prospective wet-chemistry synthesis of the top-ranked formulations is required to confirm the screening value.

\subsection{What does the deep model contribute?}

The results show a structural rather than a naive numerical advantage. Internally, SIMPLEX matches the best tree ensemble (0.79 versus 0.81, a statistical tie); prospectively, it attains the best absolute accuracy and ranking correlation of all equally tuned baselines and retains a generalisation gap approximately half that of random forest (Table 3). The mechanism is visible in the composition encoding. SIMPLEX represents composition synergy explicitly through the 15 pairwise interaction terms, so its learned response varies smoothly along the composition manifold; tree ensembles build axis-aligned, piecewise-constant partitions over raw columns and must rediscover interaction structure from data, a step that fails when a novel formulation occupies a previously empty region of the simplex.

This transfer advantage matters for the screening application. A screening model that is accurate in distribution but degrades on model-discovered formulations (random forest drops from 0.81 to 0.56) misranks precisely the candidates that an optimisation loop proposes; SIMPLEX's smaller gap (0.79 to 0.69) keeps ranking reliable exactly where screening decisions are made. The explicit interaction modality, rather than raw capacity, is the likely source of this robustness, consistent with the ablation showing multimodal fusion as the largest single contributor \cite{kim2023smallML,agrawal2016mlmaterials,isayev2017mlmaterials,ward2018mlcomposite}.

\subsection{Limitations and future work}

Several limitations warrant emphasis. (i) The prospective cohort originates from the same laboratory and instrumentation as the internal cohort (the final round of the original study's optimisation loop), so it is a model-guided prospective cohort rather than an independent-laboratory cohort; cross-laboratory validation remains an important next step \cite{liao2025nature}. (ii) The prospective cohort is small ($n = 25$), so per-metric point estimates carry uncertainty and the perfect Top-10 recovery should be confirmed on a larger cohort. (iii) Microstructural variables (cross-link density, network topology) are not recorded in the public dataset and are implicitly averaged over; incorporating microstructure descriptors is an important next step \cite{kim2023smallML,agrawal2016mlmaterials,isayev2017mlmaterials,zitnick2020materials}. (iv) The internal difference from random forest is not statistically significant, and the exact prospective ranking depends on baseline tuning; both facts are reported transparently.

(v) All conclusions are computational; prospective wet-chemistry synthesis of the top-ranked formulations is required to confirm the screening value. (vi) The prospective cohort is an optimisation-selected high-performance subset enriched in the BA-PEA region, so the Top-k metrics are conditional on this selection and are largely shared across models. Promising directions include extending evaluation to independent-laboratory cohorts, injecting microstructure descriptors as additional modalities, closing the active-learning loop by proposing the next formulations to synthesise \cite{chen2022activelearning}, and multi-target extension of the framework.

\section{Conclusion}

SIMPLEX is a dual-modality residual architecture with interaction attention and small-data regularisation that maps monomer composition to underwater adhesion strength, trained on the full experimentally explored composition-adhesion surface and validated prospectively on 25 model-discovered formulations that were never used in model selection. Under this protocol SIMPLEX matches the best-in-class tree ensemble in-distribution ($R^2$ $0.79$ vs $0.81$, a statistical tie) and attains the best prospective performance of all equally tuned baselines, with $R^2 = 0.69$, Spearman $\rho = 0.87$, perfect Top-10 recovery (1.00), and Top-20 precision 0.90, with a small internal-to-prospective generalisation gap; its interpretation analysis further distils an actionable composition rule, namely that within the explored range adhesion rises monotonically with the hydrophobic-aromatic BA$\times$PEA fraction. The framework thus supports composition screening for biohybrid interface design in data-scarce soft-materials research and development.

\section*{Conflict of Interest Statement}

The authors declare that the research was conducted in the absence of any commercial or financial relationships that could be construed as a potential conflict of interest.

\section*{Author Contributions}

T.S., X.L., and H.K.L. contributed equally to this work and are jointly listed as co-first authors. T.S. and X.L. conceived the study, designed the experimental workflow, and curated the public hydrogel adhesion dataset. H.K.L. developed the SIMPLEX architecture and implemented the training and inference pipeline. H.L. and Z.P. executed the downstream interpretability analyses, including permutation importance, attention attribution, and partial dependence profiling. X.F. supervised the project, secured funding, and wrote the manuscript. All authors reviewed the final version and approved its submission.

\section*{Funding}

This work was supported by the National Natural Science Foundation of China (No. 62505285).

\section*{Acknowledgments}

The authors gratefully acknowledge Liao et al. and the maintainers of the open hydrogel dataset for making the experimental data publicly available under the MIT licence. Constructive discussions with colleagues at the participating institutions are also appreciated.

\section*{Data Availability Statement}

The raw experimental dataset analysed in this study is publicly available from the repository \emph{sheng-hu/hydrogels} (MIT licence) accompanying Liao et al., \textit{Nature} 644, 89 to 95 (2025), DOI 10.1038/s41586-025-09269-4. The training code, the processed tensors used in this work, the per-formulation predictions underlying every reported table and figure, and the regeneration scripts are all available from the GitHub repository \url{https://github.com/SHENTongfei/HydroGelNet}. No additional materials were generated.

\bibliographystyle{Frontiers-Harvard}
\bibliography{reference_final}

\end{document}
